% Part: methods
% Chapter: proofs
% Section: reading-proofs

\documentclass[../../../include/open-logic-section]{subfiles}

\begin{document}

\olfileid[pt]{mth}{prf}{rea}
\olsection{Lendo Provas}

As provas que você encontra em livros-texto e artigos muito raramente
fornecem todos os detalhes que incluímos até agora em nossos exemplos.
Os autores frequentemente não chamam a atenção para quando distinguem
casos, quando dão uma prova indireta, ou não mencionam que usam uma
definição.  Assim, quando você lê uma prova em um livro-texto, muitas
vezes terá de preencher esses detalhes por conta própria para entender
a prova. Fazer isso também é um bom exercício para pegar o jeito dos
vários passos que você tem de dar em uma prova. Vejamos um exemplo.

\begin{prop}[Absorção]
Para todos os conjuntos $A$, $B$,
\[
A \cap (A \cup B) = A
\]
\end{prop}

\begin{proof}
Se $z \in A \cap (A \cup B)$, então $z \in A$, logo $A \cap (A \cup B)
\subseteq A$. Agora suponha $z \in A$. Então também $z \in A \cup B$, e
portanto também $z \in A \cap (A \cup B)$.
\end{proof}

A prova precedente da lei da absorção é muito condensada. Não há
menção a quaisquer definições usadas, nenhum ``temos de provar que''
antes de prová-lo, etc. Vamos desempacotá-la.  A proposição provada é
uma afirmação geral sobre quaisquer conjuntos $A$ e $B$, e quando a
prova menciona $A$ ou $B$, estes são variáveis para conjuntos
arbitrários.  As afirmações gerais que a prova estabelece são o que se
exige para provar a identidade de conjuntos, isto é, que todo
!!{element} do lado esquerdo da identidade é !!a{element} do direito e
vice-versa.

\begin{quote}
``Se $z \in A \cap (A \cup B)$, então $z \in A$, logo $A \cap (A \cup B)
  \subseteq A$.''
\end{quote}

Esta é a primeira metade da prova da identidade: ela estabelece
que, se um~$z$ arbitrário é !!a{element} do lado esquerdo, ele também é
!!a{element} do direito, isto é, $A \cap (A \cup B) \subseteq A$.
Assuma que $z \in A \cap (A \cup B)$. Como $z$ é !!a{element} da
interseção de dois conjuntos se, e somente se, é !!a{element} de ambos
os conjuntos, podemos concluir que $z \in A$ e também $z \in A \cup B$.
Em particular, $z \in A$, que é o que queríamos mostrar.  Como isso é
tudo o que precisa ser feito para a primeira metade, sabemos que o
restante da prova deve ser uma prova da segunda metade, isto é, uma
prova de que $A \subseteq A \cap (A \cup B)$.

\begin{quote}
``Agora suponha $z \in A$. Então também $z \in A \cup B$, e
portanto também $z \in A \cap (A \cup B)$.''
\end{quote}

Começamos assumindo que $z \in A$, já que estamos mostrando que, para
qualquer~$z$, se $z \in A$ então $z \in A \cap (A \cup B)$.  Para mostrar
que $z \in A \cap (A \cup B)$, temos de mostrar (pela definição de
``$\cap$'') que (i) $z \in A$ e também (ii) $z \in A \cup B$. Aqui (i) é
justamente nossa suposição, então não há mais nada a provar, e é por isso
que a prova não a menciona novamente. Para (ii), lembre-se de que $z$ é
!!a{element} de uma união de conjuntos se, e somente se, é !!a{element}
de pelo menos um desses conjuntos. Como $z \in A$, e $A \cup B$ é a união
de $A$ e $B$, este é o caso aqui. Então $z \in A \cup B$. Mostramos
ambos (i) $z \in A$ e (ii) $z \in A \cup B$, logo, pela definição de
``$\cap$,'' $z \in A \cap (A \cup B)$.  A prova não menciona essas
definições; presume-se que o leitor já as tenha internalizado.  Se você
não as internalizou, terá de voltar e relembrar quais são. Então também
terá de reconhecer por que decorre de $z \in A$ que $z \in A \cup B$, e
de $z \in A$ e $z \in A \cup B$ que $z \in A \cap (A \cup B)$.

Eis outra versão da prova acima, com tudo tornado explícito:
\begin{proof}{}
[Pela definição de $=$ para conjuntos, $A \cap (A \cup B) = A$, temos
  de mostrar (a) $A \cap (A \cup B) \subseteq A$ e (b) $A \cap (A \cup B)
  \subseteq A$. (a): Pela definição de $\subseteq$, temos de mostrar
  que se $z \in A \cap (A \cup B)$, então $z \in A$.]  Se $z \in A
\cap (A \cup B)$, então $z \in A$ [já que pela definição de $\cap$, $z
  \in A \cap (A \cup B)$ se, e somente se, $z \in A$ e $z \in A \cup B$], logo $A
\cap (A \cup B) \subseteq A$. [(b): Pela definição de $\subseteq$,
  temos de mostrar que se $z \in A$, então $z \in A \cap (A \cup B)$.] Agora
suponha [(1)] $z \in A$. Então também [(2)] $z \in A \cup B$ [já que por
  (1) $z \in A$ ou $z \in B$, o que pela definição de $\cup$ significa $z
  \in A \cup B$], e portanto também $z \in A \cap (A \cup B)$ [já que
  a definição de $\cap$ requer que $z \in A$, isto é, (1), e $z
  \in A \cup B)$, isto é, (2)].
\end{proof}

\begin{prob}
Expanda a seguinte prova de $A \cup (A \cap B) = A$, mencionando todos
os padrões de inferência usados, por que cada passo decorre de
suposições ou afirmações estabelecidas antes dele, e onde temos de
apelar para quais definições.
\begin{proof}
  Se $z \in A \cup (A \cap B)$ então $z \in A$ ou $z \in A \cap B$. Se
  $z \in A \cap B$, $z \in A$. Qualquer $z \in A$ está também $\in A \cup (A
  \cap B)$.
\end{proof}
\end{prob}

\end{document}
